\begin{table*}[t]\centering\small
\caption{Sensibilidad propia a la temperatura de partida, con PDG2025 y la misma condición neta. La referencia de transición es $156.5\pm1.5$ MeV (HotQCD, 2019); 150 MeV se conserva como inicio de la malla principal. El intervalo cercano a la transición no representa una incertidumbre completa del HRG ni un inicio de aniquilación medido.}\label{tab:initialT}
\setlength{\tabcolsep}{11pt}\renewcommand{\arraystretch}{1.14}
\begin{tabular}{rrrrr}\toprule
$T$ (MeV) & $n_B$ (${\rm cm^{-3}}$) & $n_{\bar B}$ (${\rm cm^{-3}}$) & $\Delta n_B$ (${\rm cm^{-3}}$) & $Y_\Sigma$ \\ \midrule
150 & $1.722\times10^{37}$ & $1.722\times10^{37}$ & $4.011\times10^{29}$ & 0.007465 \\
155 & $2.410\times10^{37}$ & $2.410\times10^{37}$ & $4.660\times10^{29}$ & 0.008991 \\
156.5 & $2.657\times10^{37}$ & $2.657\times10^{37}$ & $4.875\times10^{29}$ & 0.009477 \\
158 & $2.926\times10^{37}$ & $2.926\times10^{37}$ & $5.100\times10^{29}$ & 0.009976 \\
\bottomrule\end{tabular}\end{table*}
